\section{2026 年 11 月核验增补}

本节按论文首次提交时间排列。每项方法附有可核验的外部入选依据。

% METHOD: 2607.12785
\paragraph{ExtraGS：用置信加权的扩散伪视图补强内镜外推}
2026.07~\cite{2607_12785}。ExtraGS 用已有内镜高斯场在虚拟相机处渲染低质量视图，根据可见性与重叠选择值得探索的方向，再由扩散模型生成伪观测。优化同时保留真实观测，通过高斯共视和局部感知差异控制伪视图对重建的贡献。

论文在已知相机位姿和深度的 C3VDv2 结肠镜序列中评价大基线外推。扩散补全由模型先验生成，评测不能证明所补区域反映真实解剖，更不能赋予临床有效性。论文提供公开的阶段训练、虚拟视角和渲染代码。

\textbf{入选依据} \href{https://github.com/IRMVLab/ExtraGS}{公开实现}已找到与论文对应的代码文件。依据作者公开的阶段训练和渲染代码。所核对的论文版本没有列出 ByteDance Seed 作者署名。

% METHOD: 2608.07760
\paragraph{XClipGS：半空间精确积分与切面监督}
2026.08~\cite{2608_07760}。交互式切割需要处理横跨平面的高斯。按中心点硬裁会把整个原语保留或删除，容易在切面附近留下错误贡献。XClipGS 在仿射 EWA 模型下直接积分高斯落在保留半空间的部分，并用同一个可微算子训练。切割视图覆盖不同平面位置和相机距离，为原本被外层遮住的体内结构提供监督。

作者在六个 CT 和两个 MRI 体数据上比较切面边缘和错误侧漏。实验包括在同一个体内表示上替换裁剪算子，以区分算子与训练监督的作用。算子可计算新的半空间，无需再次训练，但切面视觉质量仍依赖体内表示。留出偏移位于训练区间内，这组结果主要支持对该区间内新切面的泛化。

\textbf{入选依据} 论文署名单位包括 United Imaging Intelligence（\href{https://arxiv.org/abs/2608.07760}{论文原文}）。

% METHOD: 2610.02382
\paragraph{FactorSplat：保持体渲染代理几何并编辑传递函数}
2026.10~\cite{2610_02382}。体渲染的区域传递函数把组织强度映射为颜色和不透明度。FactorSplat 保留一套高斯位置、协方差和方向相关响应，把传递函数的变化作用于视角无关颜色和不透明度。局部组织与强度查表给出物理外观改动，低秩残差弥合代理表示和原体渲染器间的差异。修改传递函数后，交互转动相机可复用更新后的高斯属性。

作者在七个 CT 与 MRI 扫描中比较留出的传递函数，报告相对于区域感知 VEG 的平均重建和编辑精度提升。逐扫描结果仍有例外，例如心脏及膝关节的编辑误差由 VEG 取得更低值。模型对所支持的传递函数族训练，更广泛的编辑需要单独验证。

\textbf{入选依据} 论文署名单位包括 United Imaging Intelligence（\href{https://arxiv.org/abs/2610.02382}{论文原文}）。
